\section{About the Topic}

\subsection{Key Terms}

\begin{guidetable}[Y{0.55}Y{1.45}]{2}
  \textbf{Term} & \textbf{Definition} \\
  \midrule
  Maritime Chokepoint & A narrow body of water that plays a strategic role as a funnel through which a large quantity of sea or ship traffic is compelled to flow. Such an area is, accordingly, susceptible to disruption and of great geopolitical significance. \\
  Exclusive Economic Zone (EEZ) & A sea area determined by UNCLOS, 200 nm from the coast, in which the state has sovereignty over resources and economic activity. \\
  Freedom of Navigation (FON) & It is the principle of freedom of passage in international waters and straits enjoyed by ships of any state without hindrance in international law. \\
  UNCLOS & United Nations Convention on the Law of the Sea, 1982 the foundation document establishing rules concerning maritime rights and jurisdictions, including territorial sea and the EEZs. \\
  Transit Passage & The right to innocent transit through an international strait without interference. It differs from the right of innocent passage through territorial waters and covers a broader range of actions. \\
  Nine-Dash Line & The vast maritime claim by the Chinese in the South China Sea that the majority was invalidated by the permanent court of arbitration in 2016, that the claim has no basis in the UNCLOS. \\
  Freedom of Navigation Operations (FONOPs) & US-backed military operations challenging the excessive territorial claims by the State. \\
  Non-State Armed Actor & An actor which is not part of or affiliated with the state, and which is not armed. \\
\end{guidetable}

\subsection{Introduction to the Topic}

Maritime choke points represent one of the most defining features of the modern global order. These strategic slivers of water, which can be either natural straits, man-made canals, or contested seaways, serve as conduits for massive share of the world’s trade to pass through.\footnote{British Institute for the Study of Iraq. (2024). Hidden in plain sight: The strategic significance and vulnerabilities of maritime chokepoints. \url{https://bisi.org.uk/reports/hidden-in-plain-sight-the-strategic-significance-and-vulnerabilities-of-maritim} e-chokepoints} The security of these chokepoints is a concern for many nations and organisations. It has become a matter of universal economic and security significance in the contemporary age.\footnote{U.S. Energy Information Administration. (2025, June 16). Amid regional conflict, the Strait of Hormuz remains critical oil chokepoint. \url{https://www.eia.gov/todayinenergy/detail.php?id=65504}} Over 80 per cent of the volume of international trade in goods is carried by sea. This number is even higher for developing countries.\footnote{UNCTAD. (2025). Review of Maritime Transport 2025: Staying the course in turbulent waters. United Nations. \url{https://unctad.org/meeting/launch-review-maritime-transport-2025}} 73 million barrels of oil are moved per day through major maritime chokepoints. This represents the vast majority of globally traded oil.\footnote{Visual Capitalist. (2025). Mapped: The world's oil chokepoints. \url{https://www.visualcapitalist.com/mapped-the-worlds-oil-chokepoints/} (citing U.S. Energy Information Administration data)} If oil cannot move through or past one of these chokepoints, it creates substantial supply delays which raises shipping costs and potentially increases world energy prices for every continent.\footnote{U.S. Energy Information Administration. (2025, June 16). Amid regional conflict, the Strait of Hormuz remains critical oil chokepoint. \url{https://www.eia.gov/todayinenergy/detail.php?id=65504}}

The five chokepoints the Council has been called upon to address today account for the flow of the vast majority of the globe's maritime trade and the bulk of energy exports that sustain the world economy.\footnote{BCG. (2024). These four chokepoints are threatening global trade. \url{https://www.bcg.com/publications/2024/these-four-chokepoints-are-threatening-global-trade}} They include: the South China Sea (and its straits), the Strait of Hormuz, the Bab el-Mandeb Strait, the Suez Canal, and the Panama Canal.

One factor that makes this challenge particularly complex for the Council is the multitude of threats involved, which some have roots in competing national sovereignty claims as between China and its neighbors in the South China Sea.\footnote{Permanent Court of Arbitration. (2016). The South China Sea Arbitration (The Republic of Philippines v. The People's Republic of China), PCA Case No. 2013-19, Award of 12 July 2016. \url{https://docs.pca-cpa.org/2016/07/PH-CN-20160712-Award.pdf}} Other threats, as in Houthi missile attacks on maritime vessels passing through the Red Sea, are posed by non-state actors.\footnote{United Nations Security Council. (2025, July 15). Adopting Resolution 2787 (2025), Security Council extends reporting mandate on Houthi attacks in the Red Sea [Press release]. \url{https://press.un.org/en/2025/sc16120.doc.htm}} And, yet others are based on the environment as the challenge of managing the low water levels of the Panama Canal or, in the operational example of the grounding of the container ship Ever Given in 2021 at the Suez Canal.\footnote{Chatham House. (2026). The maritime chokepoints that could be worse than Hormuz. The World Today. \url{https://www.chathamhouse.org/publications/the-world-today/2026-06/maritime-chokepoints-could-be-worse-hormuz}} In each of these instances lies a deeper dilemma, involving state sovereignty versus freedom of navigation and, by extension, the dilemma of collective action versus great power competition.\footnote{Security Council Report. (2025, April 30). Maritime security, May 2025 monthly forecast. \url{https://www.securitycouncilreport.org/monthly-forecast/2025-05/maritime-security-2.php}}

The tension between state sovereignty and freedom of navigation is present in every chokepoint dispute.\footnote{United Nations. (n.d.). Applying the law of the sea to protect international shipping. UN Chronicle. \url{https://www.un.org/en/un-chronicle/applying-law-sea-protect-international-shipping}} Coastal states adjacent to chokepoints attempt to assert their sovereign rights over the waters. Maritime powers insist on trade being unimpeded. The tension between maritime powers and sovereign rights continues to grow and to create more and more conflicts.\footnote{United Nations. (n.d.). Applying the law of the sea to protect international shipping. UN Chronicle. \url{https://www.un.org/en/un-chronicle/applying-law-sea-protect-international-shipping}} States have differing interpretations of UNCLOS. This is seen particularly regarding the freedom of navigation operations in the South China Sea, which has implications beyond the region.\footnote{Security Council Report. (2025, April 30). Maritime security, May 2025 monthly forecast. \url{https://www.securitycouncilreport.org/monthly-forecast/2025-05/maritime-security-2.php}}

The Security Council’s capacity to respond to chokepoint threats is constrained by the veto power of its permanent members, and how this tool has been wielded in history, along with the interconnectedness of the diplomatic parties involved in these conflicts. Economic stability relies on these chokepoints, but diplomacy is the answer to moving forward.
